%
% oracle-protocol.tex
%
% Z Specification for the z-spec Oracle Wire Protocol
% Formal model of the newline-delimited JSON contract between a Lean
% executable oracle and a property-based test harness.
%
% Type-check: fuzz -t oracle-protocol.tex
% Animate:    probcli oracle-protocol.tex -animate 20
%

\documentclass[a4paper,10pt,fleqn]{article}
\usepackage[margin=1in]{geometry}
\usepackage{fuzz}
\usepackage[colorlinks=true,linkcolor=blue,citecolor=blue,urlcolor=blue]{hyperref}

\begin{document}

\title{The Oracle Wire Protocol: A Z Specification}
\author{Formal Model of the Oracle--Harness Contract}
\date{August 2026}
\hypersetup{
  pdftitle={The Oracle Wire Protocol: A Z Specification},
  pdfauthor={Punt Labs},
  pdfsubject={Z Specification},
  pdfcreator={fuzz/probcli}
}
\maketitle

\tableofcontents
\newpage

% ============================================================
\section{Introduction}
% ============================================================

The \texttt{/z-spec:oracle} command generates two programs from a Z
specification. The first is an \emph{oracle}: a Lean executable that
holds the specification's state and answers commands about it. The
second is a \emph{driver}: a property-based test harness that sends the
oracle a random sequence of operations, sends the same sequence to the
implementation under test, and compares the two. The property being
tested is that the implementation accepts exactly what the specification
accepts.

The two programs speak newline-delimited JSON over a pipe. The oracle
emits its state before reading anything; thereafter each command it reads
yields exactly one result line. This document specifies what a result
line means.

\subsection{Why the verdict exists}

A result line carries the state \emph{and} a verdict. The verdict is not
decoration, and the reason it must be there is the subject of this
specification.

An operation whose precondition holds mutates the state. An operation
whose precondition fails leaves the state alone. But an operation whose
precondition holds may also, quite legitimately, leave the state alone
--- withdrawing nothing from an account is a lawful withdrawal of zero.
So the after-state does not determine which of the two happened. If the
line carried the state and nothing else, a rejection and a no-op success
would be byte-identical, and a driver comparing traces could not assert
the one property the harness exists to establish. An implementation that
permitted an illegal withdrawal would produce a trace indistinguishable
from one that correctly refused it, and its tests would pass.

That was a real defect in the prose statement of this protocol. Written
in Z it is an omitted output variable on two operations that share an
after-state, which is visible on inspection. Written in prose it survived
review.

\subsection{What the notation is doing}

Two conventions carry the argument.

The rejecting operation is framed with $\Xi Model$ rather than
$\Delta Model$ plus a predicate saying the state does not change. The
two are equivalent in what they permit, but they are not equivalent as
documents: with $\Xi$, ``the state is not mutated on rejection'' is a
consequence of the schema's declaration part and cannot be broken by
editing a predicate. It is a theorem, not a promise.

Every operation declares an output $ok!$, and the rejecting operation
declares a second output $reason!$. Because $ok!$ appears in the
signature of both operations, no line can be emitted without it. That is
the whole fix, stated in the one place where it cannot be forgotten.

% ============================================================
\section{Basic Types}
% ============================================================

\begin{zed}
  [PREDICATE]
\end{zed}

The names of the predicates a specification imposes. A rejection names
the predicate it violated, so that a failure report can quote it.

\texttt{PREDICATE} is deliberately a given set with no structure. The
protocol says a driver must branch on the verdict and must not parse the
reason; giving \texttt{PREDICATE} no operations is that rule stated in
mathematics rather than in a warning.

% ============================================================
\section{Free Types}
% ============================================================

\begin{zed}
  ZBOOL ::= ztrue | zfalse
\end{zed}

The verdict. \texttt{ztrue} is the JSON \texttt{true} of an accepted
operation; \texttt{zfalse} is the \texttt{false} of a rejected one.

% ============================================================
\section{Global Constants}
% ============================================================

\begin{axdef}
  maxBalance : \nat \\
  reserve : \nat \\
  reserveMaintained : PREDICATE
  \where
  maxBalance = 3 \\
  reserve = 1 \\
  reserve \leq maxBalance
\end{axdef}

\texttt{maxBalance} bounds the state and the command inputs so that
probcli can enumerate the reachable states; the protocol itself imposes
no such bound. \texttt{reserve} is the minimum balance the account must
keep. \texttt{reserveMaintained} names the precondition of
\texttt{Withdraw} --- the predicate a rejection reports.

The third constraint, $reserve \leq maxBalance$, is not arithmetic
bookkeeping. \texttt{Init} sets the balance to \texttt{maxBalance} and
the minimum to \texttt{reserve}, so it establishes the state invariant
$minimum \leq balance$ only if the two constants are ordered that way.
Discharging that obligation here, where the constants are chosen, is
cheaper than discovering it as an unreachable initial state.

% ============================================================
\section{State}
% ============================================================

The model under test is the account of the running example in
\texttt{commands/oracle.md}: a balance, and a minimum the balance may not
fall below. Only two things about it matter here. It has a precondition
worth violating, so there is something to reject; and it admits an
operation that is accepted and yet changes nothing, so acceptance and
rejection can share an after-state. Everything else about the model is
incidental to the protocol.

\begin{schema}{Model}
  balance : \nat \\
  minimum : \nat
  \where
  %
  % The invariant the rejecting case exists to protect
  %
  minimum \leq balance \\
  %
  % Bounded for ProB animation
  %
  balance \leq maxBalance
\end{schema}

\begin{schema}{Init}
  Model'
  \where
  balance' = maxBalance \\
  minimum' = reserve
\end{schema}

The oracle emits this state, with the verdict \texttt{ztrue}, before it
reads any command. The startup line is an ordinary result line; it is
described by \texttt{ReportState} below.

% ============================================================
\section{Operations}
% ============================================================

\texttt{Withdraw} is the specimen operation. Its precondition is that the
balance left afterwards still covers the minimum. The two schemas below
are the two things the oracle can do with it, and between them they cover
every input: a command is either accepted or rejected, never neither and
never both.

\subsection{AcceptWithdraw}

The precondition holds. The state changes and the verdict is
\texttt{ztrue}.

The interesting instance is $amount? = 0$. The precondition reduces to
the state invariant and therefore always holds, so a withdrawal of
nothing is always accepted --- and leaves the balance exactly where it
was.

\begin{schema}{AcceptWithdraw}
  \Delta Model \\
  amount? : \nat \\
  ok! : ZBOOL
  \where
  %
  % Bounded input for ProB animation
  %
  amount? \leq maxBalance \\
  %
  % The precondition of Withdraw
  %
  balance - amount? \geq minimum \\
  %
  balance' = balance - amount? \\
  minimum' = minimum \\
  ok! = ztrue
\end{schema}

\subsection{RejectWithdraw}

The precondition fails. The state is unchanged --- by the $\Xi$ frame,
not by a predicate --- the verdict is \texttt{zfalse}, and
\texttt{reason!} names the predicate that failed.

\begin{schema}{RejectWithdraw}
  \Xi Model \\
  amount? : \nat \\
  ok! : ZBOOL \\
  reason! : PREDICATE
  \where
  %
  % Bounded input for ProB animation
  %
  amount? \leq maxBalance \\
  %
  % The negation of the precondition of Withdraw
  %
  \lnot (balance - amount? \geq minimum) \\
  %
  ok! = zfalse \\
  reason! = reserveMaintained
\end{schema}

\subsection{ReportState}

What an observer reads off the wire: the whole state, and a verdict. This
is the startup line the oracle emits before reading any command, and it
is also the general form of an observation --- a result line reports the
state that holds after the command it answers.

\begin{schema}{ReportState}
  \Xi Model \\
  balance! : \nat \\
  minimum! : \nat \\
  ok! : ZBOOL
  \where
  balance! = balance \\
  minimum! = minimum \\
  ok! = ztrue
\end{schema}

% ============================================================
\section{The Distinguishability Property}
% ============================================================

The property the protocol exists to guarantee:

\begin{quote}
  An observer reading a result line can always tell a rejection from an
  accepted operation that happened to change nothing.
\end{quote}

The specification makes it true, and the reason is worth stating
precisely rather than asserting.

Take the state $balance = 1$, $minimum = 1$, reached from \texttt{Init}
by a single withdrawal of $2$.

\begin{itemize}
  \item \texttt{AcceptWithdraw} with $amount? = 0$ is enabled, because
    $1 - 0 \geq 1$. It leaves $balance = 1$ and emits $ok! = ztrue$.
  \item \texttt{RejectWithdraw} with $amount? = 1$ is enabled, because
    $1 - 1 \geq 1$ is false. Its $\Xi$ frame leaves $balance = 1$ and it
    emits $ok! = zfalse$.
\end{itemize}

The two lines carry the same state. Nothing recoverable from the state
distinguishes them. They differ in $ok!$ and nowhere else --- so if $ok!$
were dropped from the signatures, the two would be the same line, and no
driver could tell a specification-conforming implementation from one that
had just permitted an illegal withdrawal.

Three consequences follow from the schemas as written:

\begin{enumerate}
  \item \textbf{Every line carries a verdict.} $ok!$ is in the signature
    of \texttt{AcceptWithdraw}, \texttt{RejectWithdraw} and
    \texttt{ReportState}. There is no schema that emits a bare state.
  \item \textbf{A rejection cannot mutate.} $\Xi Model$ constrains the
    after-state to equal the before-state as part of
    \texttt{RejectWithdraw}'s declaration. No predicate in the schema can
    weaken it, and none can be forgotten.
  \item \textbf{The two cases are exhaustive and disjoint.} The
    rejecting predicate is written $\lnot P$ against the accepting
    predicate's $P$, over the same declarations, rather than as the
    equivalent $balance - amount? < minimum$. Written that way the
    argument is on the page: no input satisfies both, and every input
    satisfies one, so every command yields exactly one result line. The
    exception is an input outside $0 \upto maxBalance$, which neither
    schema admits --- that is the animation bound and not a property of
    the protocol, which accepts any $amount? : \nat$.
\end{enumerate}

% ============================================================
\section{Correspondence With the Wire}
% ============================================================

The Z above and the JSON of \texttt{commands/oracle.md} are the same
object in two notations. The correspondence:

\begin{center}
\begin{tabular}{ll}
  \textbf{Z} & \textbf{JSON} \\
  \hline
  input variables ($amount?$)     & the \texttt{args} object of a command line \\
  the after-state ($Model'$)      & the \texttt{state} object of a result line \\
  $ok! = ztrue$                   & \texttt{"ok": true} \\
  $ok! = zfalse$                  & \texttt{"ok": false} \\
  $reason!$                       & \texttt{"reason"}, present only when \texttt{ok} is false \\
  \texttt{Init} then \texttt{ReportState} & the startup line, emitted before any command is read \\
\end{tabular}
\end{center}

Serialisation of Z types into JSON values is tabulated in
\texttt{commands/oracle.md}; it is a matter of encoding and carries no
obligations that this specification does not already impose.

Two rules of the wire are outside the model, and are recorded here so
that the omission is deliberate rather than accidental. First,
\emph{one command, one line}: the framing of the pipe is a property of
the transport, not of the operations, and the disjointness argument above
is what makes it well defined. Second, \texttt{reason} is diagnostic:
\texttt{PREDICATE} is a given set precisely so that nothing can be
computed from a reason, and a driver that branches on one is relying on
something this contract does not offer.

\begin{schema}{Freeze}
  \Delta Model \\
  ok! : ZBOOL
  \where
  balance > maxBalance \\
  balance' = balance \\
  minimum' = minimum \\
  ok! = ztrue
\end{schema}

\end{document}
